\documentclass[10pt]{article}
\usepackage[a4paper,margin=0.85in,headheight=14pt]{geometry}
\usepackage{amsmath,amssymb}
\usepackage{graphicx}
\usepackage{booktabs}
\usepackage{array}
\usepackage{xcolor}
\usepackage{fancyhdr}
\usepackage{titlesec}
\usepackage{enumitem}
\usepackage{microtype}
\usepackage[hidelinks]{hyperref}

% ---------- house style (matches Layers 11-13) ----------
\pagestyle{fancy}
\fancyhf{}
\lhead{\small\itshape Layer 14: Electronics and Potentiostat Noise Budget}
\rhead{\small Aaran Nihalani $\bullet$ Eton iGEM Dry Lab}
\cfoot{\small\thepage}
\renewcommand{\headrulewidth}{0.4pt}

\titleformat{\section}{\large\bfseries}{\Roman{section}.}{0.5em}{}[\vspace{-0.6em}\rule{\textwidth}{0.4pt}]
\titlespacing*{\section}{0pt}{1.1em}{0.5em}
\titleformat{\subsection}{\bfseries}{}{0pt}{}
\titlespacing*{\subsection}{0pt}{0.7em}{-0.1em}
\setlength{\parindent}{0pt}
\setlength{\parskip}{0.45em}

\numberwithin{equation}{section}
\renewcommand{\theequation}{\Roman{section}.\arabic{equation}}
\renewcommand{\thetable}{\arabic{table}}

\newcommand{\pA}{\,\mathrm{pA}}
\newcommand{\rtHz}{\sqrt{\mathrm{Hz}}}
\newcommand{\fAV}{\,\mathrm{fA\!\cdot\! V}}
\newcommand{\nAV}{\,\mathrm{nA\!\cdot\! V}}
\newcommand{\flag}[1]{\textbf{FLAG:} #1}

\begin{document}

% ---------- title block ----------
\begin{center}
{\LARGE\bfseries Layer 14: Electronics and Potentiostat Noise Budget}\\[0.35em]
{\large From Electrode Current to Per-Point Noise, Peak-Area Noise, and the Electronic Detection Floor}\\[0.3em]
{\small\color{gray} Eton iGEM 2026, Dry Lab \;|\; Layer 14 of 17-Layer Biosensor Pipeline \;|\;
\texttt{noise\_budget.py} $\cdot$ \texttt{sigma\_area.py} $\cdot$ \texttt{snr\_vs\_gamma.py} $\cdot$ \texttt{mains\_spectrum.py}}
\end{center}
\vspace{-0.4em}\rule{\textwidth}{0.4pt}

Layer 14 sits between the electrode and the peak-extraction code. It takes the electrode capacitance and background current
from Layers 4--6 and the SWV waveform from Layers 11--12, and returns three things: (i) the RMS current noise
$\sigma_{\Delta I}$ on every net SWV point, which Layer 12 consumes as ``$\sigma_{\text{noise}}$ from Layer 14''; (ii) the
resulting noise $\sigma_{\text{area}}$ on the fitted MB peak area, which feeds Layer 13; and (iii) a decision on whether the
limit of detection is set by the electronics or by something else. The noise budget is also a design tool: it says which
component to change if the electronic floor ever becomes the problem.

\textbf{Scope and corrections.} The reference architecture is the DStat potentiostat, as specified in the planning file
(\texttt{contents.md}, Layer 14). The ProteaSign schematic (Stakeholder Document) instead shows an AD5940 analogue front
end with a TMUX1104 multiplexer, so both are evaluated (Section~I). Checking the planning file's arithmetic in code exposed several
errors (ADC step sizes off by $10^3$, electrode capacitance off by $10^1$, an omitted noise source, and an
inconsistent worked example). They are corrected here and tabulated in Section~VIII; the qualitative conclusions of the planning
file survive, but several numbers change by factors of 2--10.

% =====================================================================
\section{Signal Chain, Conventions and Inputs}

\subsection{Hardware reference}
DStat (Dryden and Wheeler, 2015 \cite{dstat}) uses an ATxmega256A3U microcontroller (not a PIC, as the planning file states),
an external 16-bit DAC (46~$\mu$V step on a 3~V range; 76~$\mu$V on 5~V), a 24-bit sigma-delta ADC whose digital filter has
zeros at multiples of its output rate, and a transimpedance amplifier (TIA) with seven selectable feedback resistors from
100~$\Omega$ to 100~M$\Omega$. Its measured open-circuit noise levels off near 200~fA (60~Hz sample rate, no cell attached),
giving an instrument detection limit near 600~fA \cite{dstat}. The ProteaSign schematic uses an AD5940 \cite{ad5940} (16-bit SAR ADC at
800~kSPS, 1.82~V internal reference, TIA resistor 0.2--512~k$\Omega$ per the schematic) behind a TMUX1104 4:1 multiplexer
(2.5~$\Omega$ on-resistance, 1.5~pC charge injection, 3~pA leakage \cite{tmux}). The DAC contributes negligibly: a
76~$\mu$V step is 3.8\% of the 2~mV staircase step, and its maximum potential error ($\pm$38~$\mu$V) is 0.2\% of the 20~mV
peak $\sigma_E$.

\subsection{Conventions}
All noise values are RMS currents referred to the working-electrode input for \emph{one} SWV current sample, taken at the end
of a half-cycle of duration $\tau = 1/(2f_{\text{SWV}})$. A boxcar average over the half-cycle has noise bandwidth
$B = 1/(2\tau) = f_{\text{SWV}} = 60$~Hz. (The planning file uses $B = f_{\text{SWV}}/2 = 30$~Hz, which would need an
integration window longer than the half-cycle.) The net SWV current is $\Delta I_n = I_{\text{fwd},n} - I_{\text{rev},n}$; the two
samples are independent, so
\begin{equation}
\sigma_{\Delta I} = \sqrt{2}\,\sigma_{\text{sample}} = \sqrt{2}\sqrt{\textstyle\sum_k \sigma_k^{2}} .
\end{equation}
Inputs, all taken from upstream layers except where marked, are in Table~\ref{tab:inputs}.

\begin{table}[h]
\centering\small
\begin{tabular}{@{}p{4.5cm}p{5.8cm}p{5.0cm}@{}}
\toprule
Parameter & Value & Source \\
\midrule
True electrode area $A_{\text{true}}$ & 0.013 cm$^2$ & Layers 4, 6 \\
SAM double-layer capacitance $C_{dl}$ & 2 $\mu$F\,cm$^{-2}$ $\Rightarrow$ $C_{\text{cell}} = 26$ nF & Layer 6 \\
Solution resistance $R_s$ & 7.1 k$\Omega$ & Layer 5 \\
$E_{sw}$, $f_{\text{SWV}}$, $\Delta E_{\text{step}}$ & 25 mV, 60 Hz, 2 mV & Layers 5, 11, 12 \\
Capacitive background $I_{\text{cap}}$ & 70 nA (60--78 nA) & Layers 5, 6 \\
Peak $\sigma_E$, window & 20 mV; $\pm3\sigma_E$ = 61 points & Layer 12 \\
MB peak area $A_{\text{MB}}$ & 0.3 nA$\cdot$V (peak 5.98 nA) & Layer 13 \\
Fc peak area $A_{\text{Fc}}$ & 1 nA$\cdot$V & Layer 12 range 0.1--10; \flag{mid-range assumed} \\
Initial coverage $\Gamma_0$ & 2.0 pmol\,cm$^{-2}$ & Layer 15 \\
Op-amp (OPA2333) & $e_n$ = 55 nV/$\sqrt{\text{Hz}}$, $i_n$ = 100 fA/$\sqrt{\text{Hz}}$, GBW 350 kHz & TI datasheet \cite{opa} \\
ADC & 24-bit sigma-delta, 20 ENOB, $\pm$2.5 V & planning file, DStat \\
Temperature (instrument) & 298.15 K & --- \\
\bottomrule
\end{tabular}
\caption{Inputs to the noise budget.}\label{tab:inputs}
\end{table}

% =====================================================================
\section{Noise Sources}

\subsection{Johnson noise in $R_f$}
\begin{equation}
i_J = \sqrt{4k_BT/R_f}
\end{equation}
gives 0.406, 0.128 and 0.041 pA/$\sqrt{\text{Hz}}$ for $R_f$ = 100~k$\Omega$, 1~M$\Omega$ and 10~M$\Omega$ at 298~K (the planning file's values
are confirmed). Over $B = 60$~Hz at 1~M$\Omega$ this is 0.99~pA RMS.

\subsection{Op-amp noise}
Two terms matter. The first is the input current noise: the OPA2333 datasheet gives $i_n = 100$~fA/$\sqrt{\text{Hz}}$ at 10~Hz \cite{opa},
5--20$\times$ larger than the 5--20~fA/$\sqrt{\text{Hz}}$ assumed in the planning file (zero-drift amplifiers trade offset and 1/$f$
performance for higher current noise). Over 60~Hz this is 0.77~pA, comparable to Johnson noise at 1~M$\Omega$, not negligible.

The second term is absent from the planning file. The WE is held at virtual ground, so the op-amp's input voltage noise $e_n$ appears across the
electrode impedance $\approx 1/(2\pi fC_{\text{cell}})$ (valid for $f \ll 1/(2\pi R_sC_{\text{cell}}) \approx 860$~Hz) and drives a
current $i(f) = 2\pi f\,C_{\text{cell}}\,e_n$ in the measurement. Integrating over the noise bandwidth,
\begin{equation}
i_{e_nC,\text{rms}} = 2\pi\,C_{\text{cell}}\,e_n\sqrt{B^{3}/3} = 2.41\pA ,
\end{equation}
using $C_{\text{cell}} = 26$~nF and $e_n = 55$~nV/$\sqrt{\text{Hz}}$ (1~kHz datasheet value). This is the \emph{largest} single term in the budget.
It grows as $B^{3/2}$, independent of $R_f$, and is the reason raising $R_f$ beyond a few M$\Omega$ buys almost nothing (Section~III).

\subsection{Shot noise of the DC background}
$i_{sh} = \sqrt{2qI_{\text{cap}}} = 0.150$~pA/$\sqrt{\text{Hz}}$ for $I_{\text{cap}} = 70$~nA, i.e.\ 1.16~pA over 60~Hz (the planning file's density is confirmed; its
summary table entry of ``0.025 at 30~Hz'' is inconsistent with it). Strictly, shot noise arises from carriers crossing a barrier; a capacitive
charging current does not have it, so this term is a conservative upper bound. The true shot noise is set by the much smaller Faradaic leakage
and O$_2$ background.

\subsection{ADC quantisation}
With input-referred step $\text{LSB} = 2V_{\text{ref}}/(R_f\,2^{\text{ENOB}})$ and noise LSB/$\sqrt{12}$: at $R_f = 100$~k$\Omega$ and $\pm2.5$~V the 20-ENOB
step is 47.7~pA (48~\emph{pA}, not the planning file's 48~nA), at 1~M$\Omega$ it is 4.8~pA, giving 1.38~pA RMS. A nominal 24-bit step at 100~k$\Omega$ is 3~pA (not 3~nA).
The planning file's lesson survives with correct numbers: at 100~k$\Omega$ quantisation dominates (14~pA RMS against $\approx$3 pA of analogue noise), so $R_f$ must be scaled up
to use the ADC range. The DStat paper reaches the same conclusion: for any ADC up to 16 bits, the step, not the analogue noise, limits resolution at every gain \cite{dstat}
(Fig.~\ref{fig:budget}C).

\subsection{1/f noise}
The OPA2333 is chopper-stabilised, so the corner is negligible. Integrating $i_n^2 f_c/f$ with $f_c = 0.1$~Hz from 0.1 to 60~Hz gives 0.08~pA. Low-frequency drift is also
suppressed by the forward--reverse differencing and removed by the baseline fit (Layer 12).

\subsection{TIA bandwidth and stability}
The planning file's $C_f$ = 1--10~pF (159~kHz at 100~k$\Omega$) ignores the 26~nF source capacitance, which destabilises a TIA. Phase margin requires
$C_f \gtrsim \sqrt{C_{\text{in}}/(2\pi R_f\,\text{GBW})} = 109$~pF at 1~M$\Omega$, so the signal bandwidth is
$1/(2\pi R_fC_f) = 1.46$~kHz, not 15.9~kHz. This is still 24$\times$ above $f_{\text{SWV}}$, so the conclusion is unchanged, but the component value must be corrected.

\subsection{Budget}
Table~\ref{tab:budget} and Fig.~\ref{fig:budget} give the base case (DStat, $R_f$ = 1~M$\Omega$, $f_{\text{SWV}}$ = 60~Hz), and the same
model for the AD5940 case. The AD5940 column uses $R_f$ = 512~k$\Omega$, 16 bits, a 1.82~V span and no PGA; \flag{its amplifier noise is a placeholder (OPA2333 values) and
the ADC span is approximate. The chip's PGA (gain 1--9) and 800~kSPS oversampling would reduce the ADC term and are not credited.}

\begin{table}[h]
\centering\small
\begin{tabular}{@{}lccc@{}}
\toprule
Source & Density (pA/$\sqrt{\text{Hz}}$) & DStat RMS/sample (pA) & AD5940 case (pA) \\
\midrule
Johnson ($R_f$) & 0.128 (1 M$\Omega$) & 0.99 & 1.39 \\
Op-amp $i_n$ & 0.100 & 0.77 & 0.77 \\
Op-amp $e_n \times C_{\text{cell}}$ & 0.539 (at 60 Hz) & \textbf{2.41} & 2.41 \\
Shot ($I_{\text{cap}}$, upper bound) & 0.150 & 1.16 & 1.16 \\
ADC quantisation & --- & 1.38 & \textbf{15.66} \\
1/$f$ & --- & 0.08 & 0.08 \\
\midrule
Total per sample & & 3.26 & 15.96 \\
\textbf{Total per net SWV point ($\times\sqrt{2}$)} & & \textbf{4.61} & 22.58 \\
\bottomrule
\end{tabular}
\caption{Noise budget, $B$ = 60 Hz, $C_{\text{cell}}$ = 26 nF. Without the $e_n C$ term (the planning-file model) the DStat total is 3.11 pA per net point.}\label{tab:budget}
\end{table}

\begin{figure}[h]
\centering
\includegraphics[width=\textwidth]{figures/fig_L14_noise_budget.pdf}
\caption{A: current noise per sample versus $R_f$ for each source. The flat terms ($e_nC$, $i_n$, shot) set a floor near 2.7 pA; the ADC term dominates below $\approx$300 k$\Omega$. B: budget at 1~M$\Omega$. C: ADC step versus analogue noise; the 16-bit step exceeds analogue noise at every $R_f$.}\label{fig:budget}
\end{figure}

% =====================================================================
\section{Choice of $R_f$}
\begin{table}[h]
\centering\small
\begin{tabular}{@{}rccccc@{}}
\toprule
$R_f$ & $\sigma_{\Delta I}$ (pA) & Input range & Stable $C_f$ & TIA BW & Dominant term \\
\midrule
100 k$\Omega$ & 20.4 & $\pm$25 $\mu$A & 344 pF & 4.6 kHz & ADC quantisation \\
1 M$\Omega$ & 4.61 & $\pm$2.5 $\mu$A & 109 pF & 1.46 kHz & $e_n C$ \\
10 M$\Omega$ & 3.97 & $\pm$0.25 $\mu$A & 34 pF & 463 Hz & $e_n C$ \\
\bottomrule
\end{tabular}
\caption{$R_f$ trade-off (DStat, 20 ENOB).}\label{tab:rf}
\end{table}
Going from 100~k$\Omega$ to 1~M$\Omega$ cuts the net noise 4.4$\times$; going on to 10~M$\Omega$ cuts it only a further 14\%, because the
$R_f$-independent terms dominate, and it would clip if the MB peak is as large as Layer 11's 0.27~$\mu$A ($R_{f,\max} = 2.5~\text{V}/0.27~\mu\text{A} = 9.3$~M$\Omega$). The planning file's recommendation of \textbf{$R_f$ = 1~M$\Omega$} is therefore confirmed,
now on the basis of range rather than noise: it is the largest decade that safely holds the biggest plausible peak. If Layer 13's smaller peak is right, autoranging to 10~M$\Omega$ is possible but unnecessary.

% =====================================================================
\section{Multiplexer Settling and Charge Injection}
When the multiplexer connects WE$_i$, the electrode charges through $R_{\text{mux}} + R_s$:
\begin{equation}
\tau_{\text{settle}} = (R_{\text{mux}} + R_s)\,C_{\text{cell}} = (2.5 + 7100)\,\Omega \times 26~\text{nF} = 185~\mu\text{s}.
\end{equation}
The planning file's 14~$\mu$s used $C = 2$~nF, ten times too small; the wait rule ``10$\tau$ $\approx$ 140~$\mu$s'' becomes 1.85~ms. The
schematic leaves idle electrodes floating, so on reconnection the electrode may sit several hundred mV from $E_{\text{start}}$. For a
0.3~V step the initial current is $0.3/7102 = 42~\mu$A, which saturates the 2.5~$\mu$A range at 1~M$\Omega$, and it takes $\ln(42~\mu\text{A}/1~\text{pA}) = 17.6\tau = 3.2$~ms
to decay below 1~pA (half-cycle at 60~Hz: 8.3~ms). Charge injection is negligible for the TMUX1104: 1.5~pC on 26~nF is a 58~$\mu$V step (2.9\% of one 2~mV step). The planning file's ADG1606 value (10--100~pC) would give up to 3.8~mV, larger than a staircase step.
(The TMUX1104 is the part in the schematic; ADG1606 is a 16:1 device.)

\textbf{Firmware requirement (replaces ``$>$100~$\mu$s''):} wait $\ge 20\tau \approx 3.7$~ms after every mux switch before the first sample, \emph{and} hold at $E_{\text{start}}$ for
$5\tau_{\text{relax}}$ (Layer 11, Sec.~IX). The second condition ($\tau_{\text{relax}} = 1/k_{\text{tot}} \sim 10$--100~ms for $k^0 \sim 10$--100~s$^{-1}$) dominates the wait. A 206-step MB-window scan at 60~Hz takes 3.4~s per electrode.
\flag{TIA recovery time from the initial saturation is not modelled and should be measured.}

% =====================================================================
\section{Mains Interference}

\subsection{Amplitude}
Mains couples in as a displacement current through a stray capacitance $C_c$ from the wiring to the high-impedance WE loop, reduced by a shield of effectiveness SE:
\begin{equation}
i_{pk} = \frac{2\pi f_m\,C_c\,\sqrt{2}\,V_{\text{rms}}}{10^{\,\text{SE}/20}} .
\end{equation}
For a UK grid (230~V, 50~Hz), $C_c$ = 1~pF gives 102~nA, reproducing the planning file's ``1--100~nA'' unshielded range. \flag{$C_c$ and SE are scenario assumptions, not measurements;
they must be measured by shorting the WE input and recording with the cell disconnected.} (The planning file quotes the shielded level as ``$<$0.1~nA'' in the text but ``$<$0.1~pA/$\sqrt{\text{Hz}}$'' in its table; a tone is not a noise density, so amplitude in nA is the right description.)

\subsection{Aliasing into the SWV waveform}
Mains is not white noise here. The net current is the difference of two samples spaced $\tau$ apart, once per staircase step. A tone $A\cos(2\pi f_mt + \phi)$ therefore appears as
$\Delta I_n = G\,A\sin(2\pi\nu n + \phi')$ with
\begin{equation}
G = 2\left|\sin(\pi f_m\tau)\right|\;\left|\mathrm{sinc}(f_mT_{\text{int}})\right|, \qquad
\nu = (f_m/f_{\text{SWV}}) \bmod 1 ,
\end{equation}
where $T_{\text{int}}$ is the averaging window of each sample and $\mathrm{sinc}\,x = \sin\pi x/\pi x$. The tone becomes a ripple in $E$-space of period $\Delta E_{\text{step}}/\nu$, which
the Savitzky--Golay filter, cubic baseline and Gaussian fit turn into an area error. For $T_{\text{int}} = \tau$, $G = 2\sin^2(\pi f_m\tau)/(\pi f_m\tau)$.
Table~\ref{tab:transfer} lists the gain and beat for a 50~Hz grid.

\begin{table}[h]
\centering\small
\begin{tabular}{@{}rcccccc@{}}
\toprule
$f_{\text{SWV}}$ (Hz) & 25 & 30 & 45 & 60 & 80 & 100 \\
\midrule
Gain $G$ & 0.000 & 0.19 & 1.11 & 1.43 & 1.41 & 1.27 \\
Alias (Hz) & 0 & 10 & 5 & 10 & 30 & 50 \\
Beat period in $E$ (mV) & DC & 6.0 & 18 & 12 & 5.3 & 4.0 \\
\bottomrule
\end{tabular}
\caption{Transfer of a 50 Hz tone into the net SWV current ($T_{\text{int}}=\tau$). On a 60 Hz grid the null is at $f_{\text{SWV}} = 30$ Hz.}\label{tab:transfer}
\end{table}

\subsection{Effect on the peak area}
Monte Carlo through the Layer 12 chain (random phase, $\pm$0.2~Hz grid detuning, tone amplitude in the linear regime) gives an area error of \textbf{0.98\% per nA of pickup} at
$f_{\text{SWV}}$ = 60~Hz on a 50~Hz grid (linear to within 3\% from 10 to 300~pA). Table~\ref{tab:mains} converts the scenarios.

\begin{table}[h]
\centering\small
\begin{tabular}{@{}lcccc@{}}
\toprule
Scenario & Pickup & \% of MB peak & $\sigma_A/A$ at 60 Hz & $\sigma_A/A$ at 45 Hz \\
\midrule
Unshielded, 1 pF & 102 nA & 1710 & $>$100\% (unrecoverable) & $>$100\% \\
Unshielded, 0.1 pF & 10.2 nA & 171 & $>$100\% & $>$100\% \\
Grounded foil, 30 dB & 0.323 nA & 5.4 & 0.30\% & 2.4\% \\
Closed cage, 60 dB & 10.2 pA & 0.17 & 0.010\% & 0.085\% \\
\midrule
\multicolumn{5}{@{}l}{Electronic floor (Section VI): $\sigma_A/A$ = 0.037\%; Layer 13 budget: 3\%} \\
\bottomrule
\end{tabular}
\caption{Mains-induced area error, 50 Hz grid.}\label{tab:mains}
\end{table}

\begin{figure}[h]
\centering
\includegraphics[width=\textwidth]{figures/fig_L14_mains.pdf}
\caption{A: pickup amplitude by scenario. B: 1~nA of 50~Hz pickup on the MB peak at three SWV frequencies (25~Hz is synchronous and leaves the peak clean). C: area error per nA versus $f_{\text{SWV}}$ for 50 and 60~Hz grids; the response is erratic, with deep nulls only at $f_m/2$.}\label{fig:mains}
\end{figure}

\subsection{Design consequences}
\begin{enumerate}[leftmargin=1.5em,itemsep=0pt,topsep=0pt]
\item \textbf{Shielding is non-negotiable.} Keeping the area error below 3\% at 60~Hz needs $\approx$31~dB, below 1\% needs $\approx$40~dB, and pushing it below the electronic floor needs $\approx$69~dB from the 1~pF open-bench case. A closed grounded cage (60~dB assumed) is sufficient.
\item \textbf{Merely detuning $f_{\text{SWV}}$ from the mains frequency does not help} (the planning file's option 2). It moves the ripple in $E$-space; on a 50~Hz grid, 45~Hz is about 8$\times$ \emph{worse} than 60~Hz because the 18~mV beat is close to the peak width.
\item \textbf{Synchronous integration works.} If each half-cycle spans exactly one mains period ($\tau = 1/f_m$, i.e.\ $f_{\text{SWV}} = f_m/2$ = 25~Hz UK, 30~Hz US), both the differencing and boxcar factors vanish and the tone and its harmonics are nulled. This needs the mux/ADC timing to be locked to $\tau = 20$~ms and is available only if $k^0_{ET}$ (Layers 10--11) puts $f_{\text{opt}}$ at or below this range. Layer 11 quotes $f_{\text{opt}} \approx k^0/2 \approx 15$~Hz for $k^0 = 30$~s$^{-1}$, which is consistent.
\item The DStat sigma-delta filter offers the same effect through its zeros at multiples of the output rate \cite{dstat}; the AD5940's sinc$^3$/sinc$^2$ filter should be configured equivalently (\flag{to be checked against the datasheet}).
\end{enumerate}

% =====================================================================
\section{Propagation to the Peak Area}
Layer 12 fits $\Delta I^{\text{corr}}(E) = A_{\text{peak}}\exp[-(E-E_0)^2/2\sigma_E^2]$ over $\pm3\sigma_E$ (61 points at 2~mV), with area $A_{\text{area}} = A_{\text{peak}}\sigma_E\sqrt{2\pi}$
and variance from its Eq.~(IV.3). For white noise $\sigma_{\Delta I}$ on each point every term of that variance scales as $\sigma_{\Delta I}^2$ and is independent of $A_{\text{peak}}$, so $\sigma_{\text{area}}$ does not depend on $\Gamma$.
Five estimates for $A_{\text{MB}} = 0.3$~nA$\cdot$V (peak 5.98~nA) and $\sigma_{\Delta I} = 4.61$~pA are in Table~\ref{tab:sigarea}.

\begin{table}[h]
\centering\small
\begin{tabular}{@{}lccc@{}}
\toprule
Estimator & $\sigma_{\text{area}}$ (fA$\cdot$V) & SNR & $\sigma_A/A$ (\%) \\
\midrule
Direct integration $\sigma_{\Delta I}\sqrt{N}\,\Delta E$ (planning file) & 72.1 & 4160 & 0.0240 \\
Free-shape Gaussian fit, Cram\'er--Rao (Layer 12 Eq.\ IV.3) & 67.3 & 4450 & 0.0224 \\
Fixed-shape fit ($E_0$, $\sigma_E$ locked), Cram\'er--Rao & 55.0 & 5460 & 0.0183 \\
\textbf{Full Layer 12 pipeline, Monte Carlo (SG $\to$ baseline $\to$ fit)} & \textbf{112.5} & \textbf{2670} & \textbf{0.0375} \\
Windowed sum after baseline subtraction, Monte Carlo & 118.7 & 2530 & 0.0396 \\
\bottomrule
\end{tabular}
\caption{Peak-area noise from electronics alone (3000 Monte Carlo trials). The full-pipeline value is adopted as the Layer 14 output.}\label{tab:sigarea}
\end{table}

The full pipeline is 1.7$\times$ worse than the Cram\'er--Rao bound, which counts neither the variance of the fitted cubic baseline nor the correlation introduced by smoothing;
this is why the Monte Carlo number is adopted. The amplitude--width correlation is $-0.58$, confirming Layer 12's remark that the covariance term must not be dropped. Three further points follow.
(i) The Gaussian fit is only 5\% quieter than a plain windowed sum, so the planning file's explanation that the fit ``extracts more information than a peak-height measurement'' is not a noise advantage of the fit.
The larger gain is from locking the peak shape (2$\times$ against the pipeline, though this gap mixes baseline and smoothing effects). (ii) $\sigma_{\text{area}} \propto \sqrt{\Delta E_{\text{step}}}$
(Fig.~\ref{fig:sigarea}C): a 1~mV step lowers it by 30\% at the price of doubling the scan time. (iii) The pipeline has a $-0.27$\% smoothing bias, common to $R(T)$ and $R(0)$ and therefore cancelled in $\rho$.
The Fc peak (1~nA$\cdot$V, 121 points) has a relative electronic noise of only 0.016\%. Layer 12's ``signal below LOD'' flag corresponds to $A_{\text{peak}} < 3\sigma_{\Delta I}/\sqrt{N_{\text{pts}}} = 1.8$~pA here, three orders of magnitude below the peak.
Layer 12's quoted Savitzky--Golay \cite{sg} sum $\sum c_k^2 = 0.233$ evaluates to 0.2075 (noise factor 0.455, not 0.48).

\begin{figure}[h]
\centering
\includegraphics[width=\textwidth]{figures/fig_L14_sigma_area.pdf}
\caption{A: Monte Carlo scatter of the fitted area against the Cram\'er--Rao width. B: estimator comparison. C: dependence on the staircase step.}\label{fig:sigarea}
\end{figure}

% =====================================================================
\section{SNR versus Probe Coverage and the Electronic Floor}
In the reversible surface-confined limit, integrating Layer 11 Eq.~(V.5) over the sech$^2$ peak gives
\begin{equation}
A_{\text{MB}}(\Gamma) = 2\,n_eF\,A_{\text{true}}\,\Gamma\,E_{sw}f_{\text{SWV}}\,f_{\text{kin}},
\end{equation}
so SNR $= A_{\text{MB}}/\sigma_{\text{area}}$ is linear in $\Gamma$. The kinetic factor $f_{\text{kin}}$ depends on $k^0_{ET}$, which is unmeasured, so it is inferred by requiring
$A_{\text{MB}}(\Gamma_0 = 2~\text{pmol\,cm}^{-2}) = 0.3$~nA$\cdot$V, giving $f_{\text{kin}} = 0.020$ (the reversible limit, 15~nA$\cdot$V, would give a 50$\times$ higher SNR).

Propagating electronic area noise through Layer 13 (Eqs.~I.4, II.4, III.5, with $n_0 = 3$ baseline scans and an intact control electrode) gives an electronic
$\sigma_{F_k}(F_k = 0)$ of \textbf{0.065\%} for one electrode pair and 0.033\% for four. The corresponding electronic $F_{\text{LOD}} = 3\sigma_{F_k}$ is 0.20\%, roughly 100$\times$ below Layer 13's
$F_{\min}$ = 20\% (Fig.~\ref{fig:snr}B). Unlike Layer 13's $(1-F_k)\times6.8\%$, the electronic $\sigma_{F_k}$ stays flat near 0.04\% out to $F_k = 0.99$, because the shrinking
peak is measured against constant absolute noise. Layer 13's $F_{\max}$ = 95\% is therefore not electronic either: the remaining probe would still have SNR $\approx$ 3 at $F_k$ = 0.9989.

\begin{table}[h]
\centering\small
\begin{tabular}{@{}lccccc@{}}
\toprule
Case & SNR at $\Gamma_0$ & $\Gamma_{\text{det}}$ (SNR=3) & $\Gamma$ ($\sigma_A/A$=3\%) & $\Gamma_{0,\min}$ ($\sigma_F<1.7\%$) & Margin \\
\midrule
DStat, full pipeline & 2670 & 0.0022 & 0.025 & \textbf{0.072} & 28$\times$ \\
DStat, fixed shape & 5460 & 0.0011 & 0.012 & 0.035 & 57$\times$ \\
AD5940 case, full pipeline & 545 & 0.011 & 0.122 & 0.352 & 5.7$\times$ \\
\bottomrule
\end{tabular}
\caption{Coverage thresholds in pmol\,cm$^{-2}$. $\Gamma_{0,\min}$ is the starting coverage at which the electronic part of $\sigma_{F_k}(0)$ reaches 1.7\% (the requirement for $F_{\text{LOD}}$ = 5\%).}\label{tab:gamma}
\end{table}

\begin{figure}[h]
\centering
\includegraphics[width=\textwidth]{figures/fig_L14_snr_vs_gamma.pdf}
\caption{A: SNR versus $\Gamma$ for the DStat (pipeline and fixed shape), the AD5940 case and the reversible limit. B: electronic $\sigma_{F_k}$ versus $F_k$ at $\Gamma_0$ = 2 pmol\,cm$^{-2}$ against Layer 13's model. C: minimum starting coverage for a target electronic $\sigma_{F_k}(0)$.}\label{fig:snr}
\end{figure}

\subsection{Is the LOD electronics-limited?}
No. Even the pessimistic AD5940 case (16-bit, no PGA, no oversampling) has SNR $\approx$ 550 at $\Gamma_0$, well above the 33 needed for 3\%. The electronic floor (0.04\%) is
80$\times$ below Layer 13's assumed 3\% per-scan area precision. With mains shielded, the mains term adds only 0.01\% (cage) to 0.3\% (partial foil). The 3\% therefore has to come from effects
that Layer 14 does not model: baseline-model error, peak drift out of the fit window, differential fouling, and incomplete redox re-equilibration between scans are the candidates,
but this layer does not establish which. The consequence for design is that improving the electronics (higher-resolution ADC, lower-noise op-amp) will not lower the LOD; the levers are upstream
(Layers 1, 9, 12) or a measurement of the actual scan-to-scan reproducibility. This also agrees with the planning file's own conclusion, but its figure ($F_{\text{LOD}} \approx 0.02$\%) is an electronics-only floor and
should not be quoted as the device LOD; the planning file's citation of a ``0.6\% LOD in Layer 13'' does not match Layer 13, which gives 20\%.

% =====================================================================
\section{Reconciliation with the Planning File}
\begin{table}[h]
\centering\footnotesize
\begin{tabular}{@{}p{3.0cm}p{3.5cm}p{3.9cm}p{4.7cm}@{}}
\toprule
Quantity & \texttt{contents.md} & This report & Reason \\
\midrule
ADC step (24 bit, $\pm$2.5 V, 100 k$\Omega$) & 3 nA; 48 nA at 20 ENOB & 3.0 pA; 47.7 pA & $10^3$ arithmetic slip \\
Op-amp current noise & 5--20 fA/$\sqrt{\text{Hz}}$ & 100 fA/$\sqrt{\text{Hz}}$ & OPA2333 datasheet \\
$e_n C_{\text{cell}}$ term & absent & 2.41 pA (dominant) & 26 nF source capacitance \\
TIA $C_f$, bandwidth & 1--10 pF; 15.9 kHz & $\approx$109 pF; 1.46 kHz & source-capacitance stability \\
$C_{dl}A$ for settling & 2 nF & 26 nF & 2 $\mu$F\,cm$^{-2}$ $\times$ 0.013 cm$^2$ \\
$\tau_{\text{settle}}$; wait & 14 $\mu$s; $>$100 $\mu$s & 185 $\mu$s; $\ge$3.7 ms + $5\tau_{\text{relax}}$ & follows from $C$; floating-electrode step \\
Charge injection & 10--100 pC (ADG1606) & 1.5 pC (TMUX1104) & part in schematic \\
Noise per point (14.8) & 2.3 pA & 3.26 pA (sample); 4.61 pA (net) & 2.3 pA used 100 k$\Omega$ Johnson; net point $\times\sqrt2$; $B$ = 60 Hz \\
$\sigma_{\text{area}}$ & 26.5 fA$\cdot$V (table: 2--5) & 67 (CR); 112 (pipeline) & table and text disagree by 10$\times$ \\
SNR at $A_{\text{MB}}$ & 30,600 (0.81 nA$\cdot$V) & 2670 (0.3 nA$\cdot$V, pipeline); 12,000 at 0.81 nA$\cdot$V (CR) & noise and amplitude differences \\
Mains shielded level & $<$0.1 nA (text) vs pA/$\sqrt{\text{Hz}}$ (table) & 10 pA tone & tone, not density \\
Detune $f_{\text{SWV}}$ to avoid mains & suggested 45 Hz & 8$\times$ worse than 60 Hz & alias beat near peak width \\
Microcontroller & PIC or similar & ATxmega256A3U & DStat paper \\
\bottomrule
\end{tabular}
\caption{Corrections to the planning file.}\label{tab:recon}
\end{table}

% =====================================================================
\section{Key Outputs}
\begin{center}\small
\begin{tabular}{@{}p{5.3cm}p{5.6cm}p{3.6cm}@{}}
\toprule
Output & Value & Feeds \\
\midrule
$\sigma_{\Delta I}$ per net SWV point & 4.61 pA (DStat, 1 M$\Omega$, 60 Hz); 22.6 pA (AD5940 case) & Layer 12 ($\sigma_{\text{noise}}$) \\
$\sigma_{\text{area}}$ (MB, full pipeline) & 112 fA$\cdot$V = 0.037\% of 0.3 nA$\cdot$V & Layers 12, 13 \\
Electronic $\sigma_{F_k}(0)$ & 0.065\% (1 pair), 0.033\% (4 pairs) & Layers 13, 15 \\
SNR at $\Gamma_0$ & 2670 (DStat); 545 (AD5940 case) & LOD is not electronics-limited \\
Minimum $\Gamma_0$ for electronic $\sigma_F < 1.7\%$ & 0.072 pmol\,cm$^{-2}$ (28$\times$ margin) & Layers 7, 9 \\
$R_f$, $C_f$ & 1 M$\Omega$, $\approx$110 pF & Hardware spec \\
Mux wait & $\ge 3.7$ ms + hold $5\tau_{\text{relax}}$ & Firmware spec \\
Shielding & Closed grounded cage; $\approx$40 dB for $<$1\% at 60 Hz & Hardware spec \\
$f_{\text{SWV}}$ & Avoid 45 Hz; 25 Hz (UK) synchronous if $k^0$ allows & Layers 10, 11 \\
\bottomrule
\end{tabular}
\end{center}

% =====================================================================
\section{Limitations}
\begin{enumerate}[leftmargin=1.5em,itemsep=1pt,topsep=0pt]
\item \textbf{Hardware is not fixed.} The ProteaSign schematic shows an AD5940, but the reference model is DStat. The AD5940 case uses placeholder amplifier noise, an approximate ADC span, and ignores its PGA and oversampling (pessimistic).
\item \textbf{$e_n$ is treated as white} at its 1~kHz value; chopper ripple and the op-amp's noise gain at higher frequencies are not modelled.
\item \textbf{Shot noise is an upper bound} because the background is capacitive charging current.
\item \textbf{Signal amplitude is uncertain by up to two orders of magnitude.} Layer 13 uses $A_{\text{MB}}$ = 0.3~nA$\cdot$V; the planning file 0.81; Layer 11's Eq.~(V.5) as written evaluates to 10.8~$\mu$A at $\Gamma = 2\times10^{-10}$ (108~nA at $2\times10^{-12}$~mol\,cm$^{-2}$) against its quoted 0.27~$\mu$A.
The SNR spans $4\times10^3$ to $2\times10^5$ over these choices, so the conclusion is robust, but $f_{\text{kin}}$ and $A_{\text{Fc}}$ need measurement.
\item \textbf{Mains scenarios are assumed}, not measured: single sinusoid, no harmonics or magnetic pickup, idealised boxcar, 0.2~Hz detuning.
\item \textbf{White, independent noise} is assumed between steps. Correlated drift, baseline-model error and other systematics are excluded, so the electronic $\sigma_F$ is a floor, not a prediction of the achieved $\sigma_F$.
\item The electronic $\sigma_F$ assumes an intact control electrode ($\rho_c = 1$), $n_0 = 3$ baseline scans and no background cleavage.
\item Floating-electrode reconnection saturates the TIA (42~$\mu$A against a 2.5~$\mu$A range); recovery time is unmodelled.
\item Layer 11 Sec.~VI's residual capacitive current (2.4~pA) appears to be a unit slip for 2.4~nA when its own inputs are used; worth checking because this feeds the baseline.
\end{enumerate}

% =====================================================================
\begin{thebibliography}{9}\small
\bibitem{dstat} Dryden MDM, Wheeler AR (2015). DStat: a versatile, open-source potentiostat for electroanalysis and integration. \emph{PLoS ONE} 10(10):e0140349. DOI: 10.1371/journal.pone.0140349. -- microcontroller, 16-bit DAC, 24-bit sigma-delta ADC, TIA ranges, measured open-circuit noise, ADC-step limitation, mains rejection by filter zeros.
\bibitem{opa} Texas Instruments. OPA333/OPA2333 1.8-V micropower CMOS operational amplifiers, zero-drift series, datasheet SBOS351. -- $e_n$ = 55 nV/$\sqrt{\text{Hz}}$ (1 kHz), $i_n$ = 100 fA/$\sqrt{\text{Hz}}$ (10 Hz), GBW 350 kHz.
\bibitem{tmux} Texas Instruments. TMUX1104 5-V, low-leakage-current, 4:1 precision multiplexer, datasheet SCDS392. -- $R_{on}$ 2--2.5 $\Omega$, charge injection 1.5 pC, leakage 3 pA.
\bibitem{ad5940} Analog Devices. AD5940/AD5941 high precision, impedance and electrochemical front end, datasheet. -- 16-bit 800 kSPS SAR ADC, 1.82 V reference, TIA current range 50 pA to 3 mA. (Only headline specifications were consulted.)
\bibitem{sg} Savitzky A, Golay MJE (1964). Smoothing and differentiation of data by simplified least squares procedures. \emph{Anal Chem} 36:1627--1639. DOI: 10.1021/ac60214a047. -- as used in Layer 12; $\sum c_k^2$ recomputed numerically.
\bibitem{hh} Horowitz P, Hill W. \emph{The Art of Electronics}, Ch.~8; and Bard AJ, Faulkner LR. \emph{Electrochemical Methods}, App.~A. -- cited by the planning file; \textbf{not consulted directly here}. The source-capacitance noise-gain term and the Schottky formula are standard results used from first principles.
\end{thebibliography}

\small Web-accessed sources (DStat, datasheets) should be saved as PDFs in \texttt{electrochemical\_device/papers/layer\_14/}, following the convention of earlier layers.

\end{document}
